\chapter{Soluciones a ejercicios selectos}
\label{ap:soluciones}

Soluciones a los ejercicios marcados con \starred{} en cada
capítulo. Los ejercicios sin \starred{} quedan abiertos
deliberadamente: la idea es que el lector se trabe un poco antes
de pedir ayuda.

\section*{Capítulo~\ref{cap:threat-hunting}: Threat hunting}

\begin{solucion}{1}
``Un usuario que se loguea en un host, desde donde se loguea en
otro host, desde donde se conecta a una IP nunca vista, y desde
esa IP hay tráfico saliente a un dominio externo.'' Tipos:
\texttt{User}, \texttt{Host}, \texttt{Host}, \texttt{IP},
\texttt{Domain}. Relaciones: \texttt{Auth}, \texttt{Auth},
\texttt{Connect}, \texttt{Beacon} (o similar).
\end{solucion}

\begin{solucion}{3}
\emph{Nested-loop} naive: \(20\,M \times 20\,M \times 20\,M = 8
\times 10^{21}\) comparaciones. Inviable. Con hash join e índices
en RAM: \(O(20\,M)\) por join, secuencial \(\Rightarrow\) tres
joins \(\sim 60\,M\) operaciones más cardinalidad intermedia. La
cardinalidad es lo que mata: en peor caso, \(20M^{1.5}\) o más,
varios GB de filas intermedias.
\end{solucion}

\section*{Capítulo~\ref{cap:grafos-temporales}: Grafos temporales}

\begin{solucion}{1}
Walks (sin restricción temporal):
\((u_1, h_1, ip_1)\) ts \((1, 4)\) y \((1, 9)\) (dos: una por cada
arista \(h_1 \to ip_1\)),
\((u_1, h_1, ip_2)\) ts \((1, 7)\),
\((u_2, h_1, ip_1)\) ts \((2, 4)\) y \((2, 9)\),
\((u_2, h_1, ip_2)\) ts \((2, 7)\),
\((u_2, h_2, ip_2)\) ts \((5, 6)\). Total: 7 walks.
\end{solucion}

\begin{solucion}{4}
\(\Delta t\)-conexidad NO es simétrica (la arista \(u \to v\) no
implica una \(v \to u\)) ni transitiva en general (un walk
\(u \rightsquigarrow v\) y otro \(v \rightsquigarrow w\) pueden
empezar muy separados en el tiempo y juntos exceder \(\Delta t\)).
La reflexividad triviales valdría definirla; pero la asimetría y
no-transitividad ya descalifican como equivalencia.
\end{solucion}

\section*{Capítulo~\ref{cap:subgraph-matching}: Subgraph matching}

\begin{solucion}{1}
En el grafo del cap.~\ref{cap:grafos-temporales} no hay
triángulos dirigidos (cada vértice es de un tipo distinto al
adyacente). Bajo isomorfismo: 0. Bajo homomorfismo: también 0
(la estructura de \texttt{User-Host-IP} no permite ciclo).
\end{solucion}

\begin{solucion}{3}
Para una cadena de \(k\) aristas, \(\rhostar = k\). El AGM
predice salida \(\leq m^k\). El bound trivial \(O(\deg^k)\)
requiere conocer \(\deg\), que en el peor caso (densidad
constante) es \(O(\sqrt{m})\) $\Rightarrow$ \(O(m^{k/2})\).
Pero AGM da \(m^k\) \emph{multiplicativo del start count}.
Para el peor caso ``todos los vértices con grado uniforme''
ambos coinciden hasta constantes. La cota AGM es ajustada via
el ejemplo del ejercicio 3 del cap. 4.
\end{solucion}

\section*{Capítulo~\ref{cap:joins-y-grafos}: Joins y AGM}

\begin{solucion}{1}
Estrella de \(k\) brazos:
\(R_0(A, B_1) \bowtie R_1(A, B_2) \bowtie \dots\). \(A\) aparece en
todas las \(k\) hiperaristas; cubrir \(A\) requiere \(\sum x_i
\geq 1\). Cada \(B_i\) aparece sólo en \(R_i\) $\Rightarrow$
\(x_i \geq 1\). Total mínimo: \(\sum x_i = k\). Pero como
\(\sum x_i \geq 1\) ya cubre \(A\), el peor cae con \(x_i = 1\)
para cada relación. \(\rhostar = k\). En grafos reales con \(A\)
muy selectivo, el bound trivial da \(\deg(A)\); rara vez ajustado.
\end{solucion}

\begin{solucion}{3}
Construcción adversarial: tomar \(R = S = T = \{(i, i+1) : i \in
[m]\}\) (\(m\) aristas en cadena lineal). Entonces el join produce
exactamente \(m - k\) tuplas (cada path de \(k\) saltos
consecutivos). Pero el AGM bound dice \(m^k\). Para acercar:
\(R, S, T = \{(i, j) : i \in [\sqrt{m}], j \in [\sqrt{m}]\}\)
(producto cartesiano de raíz). Cada relación tiene \(m\) tuplas;
el join produce \(\sqrt{m}^k = m^{k/2}\). Estricto bajo el bound
\(m^k\). Para alcanzar \(m^k\) hace falta una construcción más
balanceada.
\end{solucion}

\section*{Capítulo~\ref{cap:sistemas}: Sistemas}

\begin{solucion}{1}
Acceso secuencial: \emph{bueno}. La cache line se reusa al cargar
la próxima instancia. Acceso desde threads distintos: \emph{malo}
(\emph{false sharing}): si el thread A escribe en el campo X de la
instancia 1 y el thread B escribe en el campo Y de la instancia 2,
ambos invalidan la misma cache line en el otro CPU, pingpong
constante. Mitigación: padding al múltiplo de cache line cuando
se sabe que las instancias se acceden desde threads distintos.
\end{solucion}

\begin{solucion}{3}
Open addressing + linear probing: array contiguo de buckets;
inserts y lookups recorren linealmente desde \(\mathrm{hash}(k)\)
hasta encontrar el slot. Locality \emph{mucho mejor} para
load-factor bajo (\(< 0.5\)): el bucket inicial y los
siguientes están en la misma cache line. Para load-factor alto, las
colisiones forzan paseos largos $\Rightarrow$ peor. La hash table
``buena'' que el motor usa internamente (\texttt{ahash},
\texttt{hashbrown}) hace open addressing por esto.
\end{solucion}

\section*{Capítulo~\ref{cap:pipeline}: Pipeline}

\begin{solucion}{1}
\texttt{Ingesta}: principal en \texttt{platform/parsers} +
\texttt{platform/sources}. \texttt{Hipótesis}: \texttt{platform/dsl}.
\texttt{Matching}: \texttt{core/graph-engine}. \texttt{Scoring}:
\texttt{core/graph-engine/scoring}. La única que vive 100\% en
\texttt{core} es \emph{matching}.
\end{solucion}

\begin{solucion}{4}
Sí, hay forma. La respuesta de \texttt{run\_hunt} incluye un campo
\texttt{cap\_saturated: bool}. El frontend lo lee y muestra el
mensaje ``Estos son los primeros 9999 (cap saturado)'' o
``Resultados completos: 9999''.
\end{solucion}

\section*{Capítulo~\ref{cap:stack}: Stack de crates}

\begin{solucion}{1}
Función IP $\to$ ASN. ¿Más de un transport? Sí (Tauri y MCP la
querrán para enriquecer reports). ¿Contrato externo? Indirecto
(consume MaxMind). ¿Algorítmico puro? No (hace I/O sobre la base
MaxMind). $\Rightarrow$ \texttt{platform/enrich} (nueva crate),
NO \texttt{core}.
\end{solucion}

\begin{solucion}{4}
\texttt{tracing} en hot path es problemático: cada evento es una
asignación + serialización. Para 10K iteraciones del DFS, eso son
10K events $\Rightarrow$ \(\sim\)1 ms de overhead. Si va en
\texttt{core}, el motor depende de \texttt{tracing} (acoplamiento).
Si va en \texttt{platform}, hay que pasar trace context al matcher.
Solución pragmática: \texttt{tracing} \emph{compile-time} (feature
flag) en \texttt{core}; default-off, on para debugging.
\end{solucion}

\section*{Capítulo~\ref{cap:entity-relation}: Entity, Relation}

\begin{solucion}{1}
Sysmon Event ID 3 (Network Connection): genera 3 entidades
(\texttt{Process}, \texttt{Host}, \texttt{IP}) y 2 relaciones
(\texttt{Process -[Connect]-> IP} con metadata
\(\{\texttt{src\_port}, \texttt{dst\_port}, \texttt{protocol}\}\);
\texttt{Host -[Hosts]-> Process}).
\end{solucion}

\begin{solucion}{4}
Esquema fijo por \texttt{EntityType}: \texttt{ProcessMetadata
\{ image, cmdline, md5, sha256 \}}, \texttt{IpMetadata \{ asn,
country \}}, etc. \emph{Ganás}: ~\(50\%\) memoria (sin overhead
de hashmap), lookup más rápido (offset directo).
\emph{Perdés}: extensibilidad sin recompilar; cualquier campo
nuevo del log queda fuera. Trade-off por dominio.
\end{solucion}

\section*{Capítulo~\ref{cap:interner}: String interner}

\begin{solucion}{1}
\(10^8\) aristas \(\times\) 2 strings/arista \(\times\) 40 B
overhead per String = \(8\) GB. Con interner: \(N \times 76\) B (la
fórmula del cap.) = \(7.6 \times 10^7\) B = 76 MB. Ratio
\(\sim 100\times\).
\end{solucion}

\begin{solucion}{3}
\texttt{Vec<(Arc<str>, StrId)>} ordenado por string. Insert:
\(O(N)\) por shift. Lookup: \(O(\log N)\) por bsearch. Conviene
sólo cuando los inserts son raros pero los lookups frecuentes
(no es el caso aquí).
\end{solucion}

\section*{Capítulo~\ref{cap:compact-relation}: CompactRelation}

\begin{solucion}{1}
L2 = 1 MB = \(1\,048\,576\) B. \(\div 32\) = \(32\,768\)
\texttt{CompactRelation}s. \(\div 200\) = \(5\,242\) \texttt{Relation}s.
Ratio: \(6.25\times\).
\end{solucion}

\begin{solucion}{3}
Si quitás \texttt{repr(C)}, el compilador puede reordenar y
\emph{compactar} el struct. Mismo tamaño total ---OK---. Lo que
se rompe: el binary layout pierde estabilidad entre versiones
del compilador. Un mmap escrito con compilador A leído con
compilador B podría producir basura. La C ABI también requiere
layout fijo. \texttt{repr(C)} es la garantía contractual.
\end{solucion}

\section*{Capítulo~\ref{cap:hunt-result}: HuntResult}

\begin{solucion}{1}
\(k = 5\) saltos $\Rightarrow$ 6 vértices por path.
Vec<String>: \(10\,000 \times 6 \times (24\,B\,\text{header} +
24\,B\,\text{bytes}) = 2.88\) MB. Vec<Arc<str>>:
\(10\,000 \times 6 \times 16\,B\) = 0.96 MB \emph{más} las strings
únicas que ya viven en el interner (no se duplican). Ahorro:
\(3\times\) y elimina el costo de alloc por path.
\end{solucion}

\begin{solucion}{3}
DoS por \emph{cardinality blowup}: el atacante envía logs con
identidades únicas (\texttt{user-12345-abc}, \texttt{user-12345-abd},
etc.). Cada string entra al interner como nueva entry; el interner
crece sin bound. \texttt{Arc<str>} no mitiga esto ---es un
problema del interner, no del HuntResult---. Mitigación:
\texttt{StringInterner} con cap absoluto (default 100M strings) y
rechazo cuando se llena.
\end{solucion}

\section*{Capítulo~\ref{cap:incremental-nlf}: NLF incremental}

\begin{solucion}{1}
\(10^7 \times 36\) B = 360 MB. \(10^8 \times 36\) B = 3.6 GB.
La segunda excede budget de 1 GB por más de 3$\times$.
\end{solucion}

\section*{Capítulo~\ref{cap:dfs-baseline}: DFS}

\begin{solucion}{1}
Paths simples de longitud 3 desde \(A\):
\(A \to B \to D\) (no extiende, \(D\) sin out-edges),
\(A \to B \to E\) (no extiende, \(E\) sin out-edges),
\(A \to C \to E\) (no extiende). Total de longitud exacta 3:
0. De longitud \(\leq 3\): 3 paths de longitud 2.
\end{solucion}

\begin{solucion}{3}
\(10^4 \cdot \deg^3 = 10^4 \cdot 100^3 = 10^{10}\) operaciones
en peor caso. A 1 ns/op: 10 segundos.
\end{solucion}

\section*{Capítulo~\ref{cap:naive-matching}: Naive matching}

\begin{solucion}{1}
\(10^4 \cdot 50^4 = 6.25 \cdot 10^{10}\) iteraciones. A 1 ns/op:
~62 segundos.
\end{solucion}

\begin{solucion}{3}
Pasar por hub vs vértice promedio: hub aporta \(10^5\)
iteraciones por visita; vértice promedio aporta 5. Si el DFS
visita el hub 100 veces: \(10^7\) iteraciones extra (+1000\%).
\end{solucion}

\section*{Capítulo~\ref{cap:prune-nlf}: Prune NLF}

\begin{solucion}{1}
\(\NLF[u_1] = [\texttt{Host}=1, \dots]\),
\(\NLF[u_2] = [\texttt{Host}=2, \dots]\) (autentica en \(h_1\) y
\(h_2\)),
\(\NLF[h_1] = [\texttt{IP}=2, \dots]\) (dos aristas a \(ip_1\) +
una a \(ip_2\) $\Rightarrow$ tres),
\(\NLF[h_2] = [\texttt{IP}=1, \dots]\),
\(\NLF[ip_i] = [\dots\,\text{todo}\,0]\). Convención: cada fila
tiene 9 columnas, casi todas 0.
\end{solucion}

\begin{solucion}{3}
Mismo tipo en los 3 hijos: \(\NLF(s, L) \geq 3\). Tres tipos
distintos \(L_1, L_2, L_3\): \(\NLF(s, L_i) \geq 1\) para cada
\(i\). El motor chequea cada uno como \texttt{required}.
\end{solucion}

\section*{Capítulo~\ref{cap:prune-khop}: Prune k-hop}

\begin{solucion}{1}
Para \texttt{MAX\_DEPTH = 2}, cada vértice tiene
\(\bits[v]\) con bits prendidos según sus tipos alcanzables en
$\leq$ 2 saltos. Para \(u_1\) (autentica en \(h_1\) y \(h_2\)):
bit \texttt{Host} (vía \(h_i\)), bit \texttt{IP} (vía \(h_i \to
ip_j\)). Bits = \texttt{0b001010} (asumiendo Host=1, IP=3).
\end{solucion}

\begin{solucion}{3}
Adversarial: hipótesis con WITHIN angosto que el
\(\mathrm{min\_first}\) sólo tracking ``primer edge'' no descarta
porque el primer edge sí cae en la ventana, pero en realidad
ningún path completo cabe. La sobre-aproximación es trade-off
intencional.
\end{solucion}

\section*{Capítulo~\ref{cap:wcoj}: WCOJ y AGM}

\begin{solucion}{1}
\(G(1000, 0.05)\): \(m = \binom{1000}{2} \cdot 0.05 \approx
25\,000\) aristas. Triángulos en expectation: \(\binom{n}{3}
p^3 = \binom{1000}{3} \cdot 0.05^3 \approx 20\,000\). AGM bound:
\(m^{1.5} \approx 4 \cdot 10^6\) ---muy por encima.
\end{solucion}

\begin{solucion}{3}
Hash join binario para \(R \bowtie S\) (sobre B): cardinality
intermedia hasta \(|R| \cdot |S|\) si todos los B coinciden. Para
\(|R| = |S| = m\): hasta \(m^2\). Sigue ahí cuando se joina con
\(T\), después se filtra ---pero el peak intermedio fue
\(m^2\). Memory blowup.
\end{solucion}

\section*{Capítulo~\ref{cap:lftj}: LFTJ}

\begin{solucion}{1}
\(K_4\)s en \(G(1000, 0.05)\): \(\binom{n}{4} p^6 = \binom{1000}{4}
\cdot 0.05^6 \approx 0.65\). Esperaba 1 \(K_4\) en average. AGM
bound: \(m^2 = 6.25 \cdot 10^8\) ---enormemente sobre-estimado.
\end{solucion}

\begin{solucion}{3}
DFS para \(K_4\) en peor caso: itera todos los starts \(\cdot\)
todos los segundo-vecinos \(\cdot\) etc. = \(O(n \cdot \deg^3)\)
= \(O(n \cdot (m/n)^3) = O(m^3 / n^2) \geq O(m^2)\) cuando
\(n \leq \sqrt{m}\) (graphs densos). El DFS bien-podado puede
mejorar la constante pero no el bound asintótico sin información
estructural extra.
\end{solucion}

\section*{Capítulo~\ref{cap:hipergrafos}: Hipergrafos}

\begin{solucion}{1}
Estrella: hiperaristas \(\{c, h_1\}, \{c, h_2\}, \{c, h_3\}\).
\(h_i\) aparece sólo en \(F_i\) $\Rightarrow$ eliminar; quedan
\(\{c\}, \{c\}, \{c\}\). Iguales $\Rightarrow$ eliminar dos por
contención. Queda \(\{c\}\); eliminar \(c\) (aislado en una
hiperarista). \(F = \emptyset\). \(\alpha\)-acíclica.
\end{solucion}

\begin{solucion}{3}
Hipergrafo \(\alpha\)-acíclico de 4 relaciones: una cadena lineal
\(R_1, R_2, R_3, R_4\) tiene 1 join tree (lineal). Una estrella
con root \(R_0\) y hojas \(R_1, R_2, R_3\): 1 join tree.
Combinaciones más complejas pueden tener varios; en general el
join tree NO es único.
\end{solucion}

\section*{Capítulo~\ref{cap:yannakakis-chain}: Yannakakis chain}

\begin{solucion}{1}
50 matches de output, \(|B_0| = 1000\), \(|B_1| = 5000\). Bottom-up:
\(B_0\) queda con sólo las rows cuyo dest aparece como origin en
\(B_1\). En peor caso \emph{ninguna} row se filtra (todas las
1000 tienen continuación). En mejor caso: \(B_0\) reducido a
\(\leq 50\) (sólo las que llevan a un match). Top-down: similar.
\end{solucion}

\begin{solucion}{3}
Si \(B_0\) queda vacío: no hay starts, el enumerator no produce
nada. Demuestra: \(B_0\) vacío post-bottom-up significa que ninguna
row tiene continuación; para que un match exista, su row inicial
debe estar en \(B_0\) y tener continuación. Contradicción.
Converse NO vale: \(B_0\) no vacío no garantiza matches; algunos
rows pueden tener continuación parcial pero fallar la cadena
completa.
\end{solucion}

\section*{Capítulo~\ref{cap:yannakakis-tree}: Yannakakis tree-walk}

\begin{solucion}{1}
\(c\) tiene 100 vecinos \texttt{Host}. Para
\texttt{c -[Auth]-> h1; c -[Auth]-> h2}: cada combinación
ordenada de 2 hosts distintos cuenta como match
$\Rightarrow$ \(100 \times 99 = 9\,900\) matches. Cota AGM:
estrella con \(\rhostar = 2\), input \(m = 100\) per relation,
output \(\leq m^2 = 10^4\). Cerca del bound.
\end{solucion}

\begin{solucion}{3}
Caso \texttt{(Some, Some)}: hipótesis donde el step actual tiene
ambos extremos ya bound. Por ejemplo, hipótesis triangular
\texttt{a -[E1]-> b; b -[E2]-> c; a -[E3]-> c}: el último step
\texttt{a -[E3]-> c} tiene \(a\) bound (de step 0) y \(c\) bound
(de step 1). Lookup determinista: ¿existe la arista? Si sí,
recurse; si no, fail.
\end{solucion}

\section*{Capítulo~\ref{cap:planner}: Planner}

\begin{solucion}{1}
\texttt{User -[Auth]-> Host -[Connect]-> IP}: shape Chain,
\(\alpha\)-acíclica $\Rightarrow$ Yannakakis chain.
\texttt{c -[Knows]-> p1; c -[Knows]-> p2; p1 -[Knows]-> p2}:
shape Graph, cíclica (triángulo) $\Rightarrow$ DFS (no aplica
LFTJ, no aplica Yannakakis).
``k4-clique'' con 6 pasos heterogéneos: name marker dispatch
$\Rightarrow$ K4Lftj (back-compat); pero structural\_dispatch
fallaría por heterogeneidad. El name marker prevalece.
\end{solucion}

\begin{solucion}{3}
Detector de estrella de \(k\) brazos: chequear que la hipótesis
shape graph tenga exactamente \(k\) cláusulas, todas con un
binding compartido como root, y los demás vértices únicos. Función
booleana: \texttt{is\_star(hyp, k) -> bool}.
\end{solucion}

\section*{Capítulo~\ref{cap:indices-auxiliares}: Diseño índices}

\begin{solucion}{1}
La sesión NO se rompe; el grafo canónico (\texttt{edge\_store})
sobrevive. Próximo hunt: \texttt{ensure\_nlf} reconstruye desde
cero. Costo: $\sim$30 segundos extra. Funcionalidad intacta.
\end{solucion}

\begin{solucion}{3}
\texttt{first\_seen[v]}: \texttt{Vec<i64>} indexado por sid.
Construye con una pasada \(O(|\E|)\) tomando min de timestamps
para cada source. Update incremental: cada
\texttt{add\_relation} hace \texttt{first\_seen[v] =
min(first\_seen[v], cr.timestamp)}. Trivial, lineal, atómico-free
single-writer.
\end{solucion}

\section*{Capítulo~\ref{cap:nlf-detalle}: NLF en detalle}

\begin{solucion}{1}
\(5 \cdot 10^7 \cdot 36\) B = 1.8 GB. Excede budget default 1 GB
$\Rightarrow$ spillea a mmap.
\end{solucion}

\begin{solucion}{3}
NLP (Neighborhood Label Profile): \texttt{[u32; 9 \(\times\) 16]}
si tenemos 16 RelTypes. = 576 B per row. \(\sim 16\times\) más
memoria. Para \(10^7\) vértices: 5.4 GB ---supera budget. Sólo
viable con dense+sparse layout o sin spill.
\end{solucion}

\section*{Capítulo~\ref{cap:khop-detalle}: k-hop en detalle}

\begin{solucion}{1}
\(|\V| = 5 \cdot 10^7\), 50\% ranked $\Rightarrow$
\(|\V_{\text{ranked}}| = 2.5 \cdot 10^7\). bits: 100 MB. rank:
200 MB. min/max\_first sparse: \(2 \cdot 2.5 \cdot 10^7 \cdot 72\)
B = 3.6 GB. Total: 3.9 GB.
\end{solucion}

\begin{solucion}{3}
Adversarial: hipótesis con WITHIN angosto pero
\(\mathrm{min/max\_first}[v][L]\) wide. El primer edge cae en la
ventana (no descartado), pero el segundo edge (que el motor
todavía no chequea con este índice) cae fuera. El DFS lo
descarta después de empezar. Trade-off intencional.
\end{solucion}

\section*{Capítulo~\ref{cap:dsl-design}: Diseñando una DSL chica}

\begin{solucion}{1}
\texttt{Process p -[Write]-> File f; f -[Connect]-> IP WITHIN 5m}.
Componentes: cadena \texttt{p -> f -> IP} (graph shape por el
\texttt{;}), ventana \texttt{WITHIN 5m}, sin predicados.
\end{solucion}

\section*{Capítulo~\ref{cap:dsl-chain}: Sintaxis chain}

\begin{solucion}{1}
\texttt{User u -[Auth]-> Host h WITHIN 10m HAVING count(distinct h) >= 3}.
Cláusulas: WITHIN para la ventana, HAVING para la cota.
\end{solucion}

\begin{solucion}{3}
\texttt{Hypothesis} con 2 \texttt{HypothesisStep}s: step 0 con
\texttt{source\_type=Process, dest\_type=Process, rel=Spawn}; step
1 con \texttt{source\_type=Process, dest\_type=File, rel=Write}.
\texttt{shape=Chain}, \texttt{window\_seconds=Some(60)},
\texttt{k\_simplicity=1}, \texttt{aggregations=vec![]}.
\end{solucion}

\section*{Capítulo~\ref{cap:dsl-graph}: Sintaxis graph}

\begin{solucion}{1}
\texttt{User u -[Spawn]-> Process p; p -[Write]-> File f; p -[Spawn]-> Process c}.
Bindings compartidos: \(p\) aparece en cláusulas 1, 2, 3 (root del
join tree); \(u\), \(f\), \(c\) sólo en una cláusula cada uno.
\end{solucion}

\begin{solucion}{3}
\texttt{User c -[Auth]-> Host h; c -[Auth]-> Host h}: dos cláusulas
\emph{idénticas}. Legal. El matcher trata ambas como restricciones
sobre los mismos bindings $\Rightarrow$ equivalente a una sola
cláusula. Output: matches sin duplicación.
\end{solucion}

\section*{Capítulo~\ref{cap:ast}: AST al matcher}

\begin{solucion}{1}
\texttt{Hypothesis\{ steps: [Step(c, Auth, h1), Step(c, Auth, h2)],
shape: Graph, ...\}}. \texttt{VarId}s: \(c \to 0\), \(h_1 \to 1\),
\(h_2 \to 2\).
\end{solucion}

\begin{solucion}{3}
API REST: \texttt{POST /hypothesis/parse} con
\texttt{\{ "dsl": "..." \}} $\to$ \texttt{Hypothesis} JSON;
\texttt{POST /run\_hunt} con
\texttt{\{ "hypothesis": ... \}} $\to$
\texttt{\{ "results": [...], "cap\_saturated": false \}}. Errores:
400 Bad Request con detalle estructurado (\texttt{\{ "error":
"unknown\_binding", "hint": "did you mean: c, h1?" \}}).
\end{solucion}

\section*{Capítulo~\ref{cap:khop-delta-hard}: Khop delta es difícil}

\begin{solucion}{1}
Hub de in-degree \(10^4\). Si la nueva arista agrega un tipo
nuevo a \(u\), todos los \(10^4\) predecesores del hub lo ven.
Cada uno propaga a sus predecesores; en grafo scale-free con
distribución similar, ronda 2: \(10^4 \cdot 10 = 10^5\). Ronda 3:
\(10^6\). Ronda 4: \(10^7\). Cascada explota.
\end{solucion}

\begin{solucion}{3}
Si \(\bits[v] \subseteq \bits[u]\): no hace falta propagar bits
(monotonicidad: ya están). Pero \(\mathrm{min/max\_first}\) puede
todavía mejorar: \(t\) podría ser timestamp más temprano que
cualquier camino actual desde \(v\). Hay que chequearlo
separadamente.
\end{solucion}

\section*{Capítulo~\ref{cap:delta-wcoj}: Delta WCOJ}

\begin{solucion}{1}
\(\Delta_u = \bits[v] \setminus \bits[u] = \texttt{0b001000}
\setminus \texttt{0b000001} = \texttt{0b001000}\) (sólo el bit IP
es nuevo). \(\bits[u] \mathrel{|}\!= \bits[v] = \texttt{0b001001}\).
\end{solucion}

\begin{solucion}{3}
Inducción + monotonicidad. Base: \(N = 0\), trivial. Paso: cada
\texttt{delta\_update} agrega exactamente los bits del walk
introducido por la nueva arista. La unión de las propagaciones
secuenciales = la propagación única que un rebuild haría.
Monotonicidad garantiza que el orden no afecta. \(\square\)
\end{solucion}

\section*{Capítulo~\ref{cap:streaming-canonical}: Streaming canonical}

\begin{solucion}{1}
\(2 \cdot 10^8 \times 32\) B = 6.4 GB ahorrados.
\end{solucion}

\begin{solucion}{3}
Tests ``correcto'': comparar resultados de
\texttt{run\_hunt} sobre streaming vs sobre edge\_store, esperar
iguales. Tests ``rápido'': bench tiempo de
\texttt{matching\_hops} sobre streaming, comparar geomean vs
edge\_store, gate \(\leq 10\%\) regresión.
\end{solucion}

\section*{Capítulo~\ref{cap:plug-and-play}: Plug-and-play}

\begin{solucion}{1}
``Cache de NLF compartido entre hipótesis similares'': cumple
default-on (cache se enciende solo). Surface UI: agregar
\texttt{cache\_hits/misses} a \texttt{PerfStatus}. Default-off:
no aplica si está bien probado. Doc: actualizar
\filepath{08-env-vars.tex}. Capabilities: no aplica (no es feature
flag de build).
\end{solucion}

\begin{solucion}{3}
Test: parsear el changelog del PR; para cada nuevo flag detectado,
chequear que (a) está en la lista default-on de \texttt{lib.rs}, (b)
está en \texttt{PerfStatus}, (c) está documentado en
\filepath{08-env-vars.tex}. Falla CI si falta cualquiera de los
tres.
\end{solucion}

\section*{Capítulo~\ref{cap:frontend-instr}: Frontend instr}

\begin{solucion}{1}
``Build: AVX2'' pildora: agregar \texttt{has\_avx2: bool} a
\texttt{PerfStatus} (vía \texttt{cfg!(target\_feature = "avx2")}).
Render: pildora gris ``Build: AVX2'' si \texttt{has\_avx2}
$\Rightarrow$ true; ocultar si false.
\end{solucion}

\begin{solucion}{3}
Polling cada 5 s: balance entre responsividad
(detectar pico antes de OOM) y costo (cada poll cuesta \(\sim 1\)
ms en CPU + redraw). Más frecuente: overhead percibido. Menos:
podemos miss un pico antes de que el usuario entre en swap.
\end{solucion}

\cleardoublepage
